\section{Further proofs and details of the theory}
\label{app:proofsmore}

This section completes Appendix~\ref{app:proofs} of the main text.

\subsection{The two ends of the hierarchy}

With the counts $c_{\mathcal C}(i,j)$ and pair fractions $\pi_{\mathcal C}(c)$ of the
Omega index (Sec.~\ref{sec:metrics}):

\begin{lemma}
\label{lem:constant}
If $c_{\mathcal C}(i,j)=c$ for every pair and $c_{\mathcal G}$ is not identically $c$,
then $\Omega=0$.
\end{lemma}

\begin{proposition}[Endpoints]
\label{prop:endpoints}
Let the ground truth give some pair count $0$ and some pair count $\ge1$.
(i) If $G$ is connected and $|V|\ge3$, then $\mathcal C_0=\{V\}$ and $\Omega(0)=0$.
(ii) The communities at $t=1$ are the classes of nodes with identical closed
neighbourhoods having at least three members; if there are none,
$\mathcal C_1=\varnothing$ and $\Omega(1)=0$.
Consequently, when both hold, whenever $\Omega(t)>0$ for some $t$ the maximum
$\Omega^\star=\max_t\Omega(t)$ is attained at an interior threshold $t^\star\in(0,1)$.
\end{proposition}
The proposition does not assert that $\Omega^\star>0$ or that $\Omega(t)$ has a single
maximum; both are empirical. On noise-free ring lattices (ii) holds when $s\ge2k+2$, and
(i) when the merges connect all rings, which random merges usually do.

\begin{proof}[Proof of Lemma~\ref{lem:constant}]
$\pi_{\mathcal C}$ is concentrated on $c$, so $\Omega_e=\pi_{\mathcal G}(c)<1$; a pair agrees
exactly when $c_{\mathcal G}(i,j)=c$, so $\Omega_u=\pi_{\mathcal G}(c)$ as well.
\end{proof}

\begin{proof}[Proof of Proposition~\ref{prop:endpoints}]
(i) At $t=0$ every wedge is admissible. Two edges sharing an endpoint form a wedge, and
the wedges built on a common edge share its endpoints, so a walk in the line graph of
$G$ is a walk in $L_0$. The line graph of a connected graph is connected, and every
node of a connected graph with $|V|\ge3$ lies on a wedge. The single community $V$
gives $c_{\mathcal C}\equiv1$, and Lemma~\ref{lem:constant} applies.
(ii) $J(u,v)=1$ iff $N[u]=N[v]$, an equivalence relation whose classes are cliques; a
triple is admissible at $t=1$ iff it lies within one class, and triples within one
class are connected in $L_1$. Without classes of size at least three, $c_{\mathcal C}\equiv0$
and Lemma~\ref{lem:constant} applies. For the final claim: the cover changes only at
the finitely many distinct triple scores, so $\Omega(t)$ takes finitely many values on
$[0,1]$ and its maximum $\Omega^\star$ is attained; under the stated hypotheses
$\Omega(0)=\Omega(1)=0$ by (i) and (ii), so since $\Omega^\star\ge\Omega(t)>0$ every
maximiser lies in $(0,1)$.
\end{proof}

\subsection{Further ring-lattice windows}

\emph{Smaller rings and multiple merges.} For a single merge between two rings, exact
computation for $k=4$ to $12$ shows that the window of Eq.~\eqref{eq:window} also holds
for ring sizes down to $s=k+4$ for $k=4$ and $6$, and $s=k+2$ for $k=8$ to $12$; smaller
rings widen it upward (for $k=4$ and $s=6$ or $7$ it is $(1/9,2/5]$), which covers the
rings of Fig.~\ref{fig:schematic}. A node formed by merging one node from each of $m$
rings (all other merges isolated) gives the window
$(1/(2k+1),\min\{k/(mk+2),\max((k-2)/(k+4),(k-1)/(mk+3))\}]$, whose upper edge tends to
$1/m$ as $k$ grows; it is non-empty if and only if $m\le2k$. These statements and
Eqs.~\eqref{eq:window} and~\eqref{eq:windowab} were checked in rational arithmetic at
and on either side of every distinct commonality value (analysis/exact\_windows.py).

\emph{Unequal degrees.} For one isolated merge between rings of even degrees
$k_a,k_b\ge4$ and sizes at least $2k_a+2$ and $2k_b+2$, the same argument applies with
the commonalities of Lemma~\ref{lem:ringJ} for unequal degrees: with $x,z$ on the ring
of degree $k_a$, $J_u(d)=(k_a+1-d)/(k_a+k_b+1+d)$ and $J(z,y')=1/(k_a+k_b+1)$. The cross
pair has
$J=1/(k_a+k_b+1)$, which gives the lower edge of Eq.~\eqref{eq:windowab}. On the ring of
degree $x$, with $y$ the degree of the other ring, the upper-edge argument runs with
$J_u(1)=x/(x+y+2)$, $J_u(2)=(x-1)/(x+y+3)$ and $J(3)=(x-2)/(x+4)$, which gives
$U(x,y)$; both rings must be recovered, which gives the minimum over $U(k_a,k_b)$ and
$U(k_b,k_a)$.

\emph{Obstruction for unequal degrees.} A shared node between a dense and a sparse
community is held more weakly by the sparse one, so the window of such a merge narrows
and moves down, with upper edge $k_b/(k_a+k_b+2)$ for $k_a\gg k_b$. It cannot overlap
the window of merges between two sparse rings, whose lower edge is $1/(2k_b+1)$, once
$k_a\ge2k_b^2-2$: beyond that degree contrast no single threshold recovers a family that
also contains merges between two sparse rings, even without noise.
\subsection{Fusion of rings through shared pairs}

\emph{Fusion of rings.} Connectivity through shared pairs is what lets a community grow
through its own open wedges, and Proposition~\ref{prop:sharedpairs} is its counterpart.
On ring lattices, three communities that pairwise share single
nodes are joined when two of the three shared-node pairs are adjacent and the third lies
within distance $k$; for $k=8$ and pairwise adjacent shared nodes two apart, the rings
are recovered only for $4/13<t\le4/9$, and for $k=4$ they are never recovered. Two rings
joined only through two shared merged nodes, with no other merge within ring distance
$k$ of them, can be fused above the lower edge $1/(2k+1)$ only when those nodes lie
within ring distance $k$ of each other on both rings, or within $k/2$ on one ring and
within $3k/2$ on the other; for $k=8$, shared nodes $2$ apart on one ring and $11$ apart
on the other fuse the rings for $t\le3/23$. Theorem~\ref{thm:window}'s isolation
hypothesis excludes both cases. Whether fusion then occurs depends on the threshold,
since the merge lowers the commonality of the shared pair: two rings with $s=50$, $k=8$
that share two nodes four positions apart are fused at small $t$ but recovered
separately at $t=0.4$, where the shared pair, with $J=1/3$, is not admitted.

\subsection{Cost}

Every commonality needed is available from the wedges themselves: for a pair $\{a,b\}$,
$|N[a]\cap N[b]|$ is the number of common neighbours of $a$ and $b$, that is of centred
wedges $a$--$c$--$b$ (a triangle contributes once to each of its three pairs), plus $2$
if $a\sim b$, so one pass over the centred wedges, with pairs indexed by hashing, counts
all of them. Our implementation finds components on the incidence graph between admitted
triples and their pairs over a grid of thresholds. Alternatively, as in sequential clique
percolation~\cite{kumpula2008sequential}, a single union--find sweep over the triples in
decreasing order of score builds the merge tree of the triple components at all
thresholds in $O(W\log W)$, from which the cover at any one threshold is read in $O(W)$
more. For heavy-tailed degree distributions the hubs dominate through $W$, as for
every triangle-based method.

\emph{Degree-ratio pruning.} In the pruning of Sec.~\ref{sec:cost}, every wedge through
an edge with $\rho<t_{\min}$ is discarded. Common neighbours reached through discarded
edges must still be counted, for example by intersecting the adjacency lists of the kept
pairs; counting only the kept wedges would underestimate $J$. This is the size filter of
exact set-similarity joins~\cite{arasu2006efficient,bayardo2007scaling}, applied to
closed neighbourhoods of sizes $d_u+1$ and $d_v+1$. The pruning applies to every rule
that tests the edges of a wedge, including our three ablations, but not to link
communities, which score a wedge by its non-adjacent pair alone, whose commonality is
not bounded by the degrees of the wedge's edges, so every wedge through a hub has to be
enumerated. Our implementation counts the common neighbours of the kept pairs from the
full adjacency matrix, and pruned and unpruned runs give identical covers on its test
graphs, which include a hub joined to many leaves, and on the full Amazon and DBLP graphs
at all 30 thresholds $t\ge0.02$ of our grid.

